\documentclass[10pt]{article}

% ---------- PACKAGES ----------
\usepackage[margin=1.5cm]{geometry}
\usepackage{amssymb, amsthm}
\usepackage{mathtools}
\usepackage{fancyhdr}
\usepackage{enumitem}

% ---------- CUSTOM COUNTERS & COMMANDS ----------
\newcounter{problemnum}
\newcounter{subpartnum}[problemnum]

\newcommand{\problem}[1]{
  \stepcounter{problemnum}
  \setcounter{subpartnum}{0}
  \vspace{0.8em}
  \noindent
  {\boldmath{\textbf{#1}}}
  \nopagebreak
  \vspace{0.4em}
  \par
}

\newcommand{\subpart}{
  \stepcounter{subpartnum}
  \vspace{0.3em}
  \noindent
  \textbf{(\alph{subpartnum})}
  \quad
}

\newenvironment{solution}{
  \vspace{0.2em}
  \noindent
  \normalsize
  \textit{Solution:}
  \nopagebreak
  \begin{list}{}{
    \setlength{\leftmargin}{2em}
    \setlength{\parsep}{0.8em}
    \setlength{\itemindent}{0em}
  }
  \item
}{
  \end{list}
  \vspace{0.3em}
}

% BOLD MATH STATEMENT COMMAND
\newcommand{\boldstatement}[1]{
  {\boldmath \underline{\textbf{#1}}}
}

% RIGHT-SIDE JUSTIFICATION COMMAND
\newcommand{\justify}[1]{
  \hfill \textit{(#1)}
}

% Alternative: For multi-line proofs with justifications
\newcommand{\justifyline}[2]{
  #1 \hfill \textit{(#2)} \\
}

% ---------- GLOBAL STYLE ----------
\setlength{\parindent}{0pt}
\setlength{\parskip}{0.5em}
\pagestyle{fancy}
\fancyhf{}
\cfoot{\thepage}
\renewcommand{\headrulewidth}{0pt}

% ---------- DOCUMENT INFO ----------
\newcommand{\courseName}{MATB24H3}
\newcommand{\assignmentNum}{1}
\newcommand{\studentName}{Your Name}
\newcommand{\studentNumber}{Your Student Number}
\newcommand{\dueDate}{Due Date}

% ---------- TITLE ----------
\title{\textbf{\courseName}}
\author{\studentName}

% ----- SHORTHAND ----------
\newcommand{\Q}{\mathbb{Q}}
\newcommand{\Z}{\mathbb{Z}}
\newcommand{\R}{\mathbb{R}}
\newcommand{\C}{\mathbb{C}}
\newcommand{\F}{\mathbb{F}}
\newcommand{\N}{\mathbb{N}}
\newcommand{\K}{\mathbb{K}}
\newcommand{\cM}{\mathcal{M}}
\newcommand{\cF}{\mathcal{F}}
\newcommand{\cR}{\mathcal{R}}
\newcommand{\cP}{\mathcal{P}}
\newcommand{\cL}{\mathcal{L}}
\newcommand{\cC}{\mathcal{C}}
\newcommand{\cB}{\mathcal{B}}

\begin{document}

\begin{center}
  \LARGE\textbf{MATB24H3}\\
  \Large\textbf{Exercise Questions}
\end{center}

\textbf{Student Name: \studentName} \hfill \textbf{Student Number: \studentNumber}
\vspace{0.5em}
\hrule

% ---------- USAGE INSTRUCTIONS ----------
% Please be responsible and don't include any solutions in this file when uploading.
%
% You can use this file to write your own solutions in a separate file.
%
% Please use the features \problems, \subpart, \begin{solution}, \justify, \boldstatement{}, etc. to write your solutions in a clear and organized manner.
%
% Please help gather all the problems for the course into this file for the good of the public.
%
% The following are shorthand for the following commands:
%
%   \mathbb{Q} => \Q
%   \mathbb{Z} => \Z
%   \mathbb{R} => \R
%   \mathbb{C} => \C
%   \mathbb{F} => \F
%   \mathbb{N} => \N
%   \mathbb{K} => \K
%   \mathcal{M} => \cM
%   \mathcal{F} => \cF
%   \mathcal{R} => \cR
%   \mathcal{P} => \cP
%   \mathcal{L} => \cL
%   \mathcal{C} => \cC
%   \mathcal{C} => \cB

% ---------- PROBLEMS ---------- 
\problem{
  2.1.1. Consider the set $X = \Q \setminus \{0\}$ and the map $\star$ is depicted by $a \star b = \frac{a}{b}$. Is $\star$ an associative binary operation?
}
\begin{solution}
  % ==== Place for solution here ====
\end{solution}

\problem{
  2.1.2. Is $(\cP(\N), \cap)$ a monoid? what is the identity element? How about $(\cP(X), \cup)$ for any arbitrary set $X$?
}
\begin{solution}
  % ==== Place for solution here ====
\end{solution}

\problem{
  2.2.1. Suppose $d$ is square-free, so there is no integer $m$ such that $m^2 = d$. Show that ($\Q(\sqrt{d})$, +, ·)
forms a field. What happens if $d = m^2$ for some $m$?
}
\begin{solution}
  % ==== Place for solution here ====
\end{solution}

\problem{
  2.2.2. Now, can we extend a finite field in the same way we extended an infinite field? If so, give an example.
}
\begin{solution}
  % ==== Place for solution here ====
\end{solution}

\problem{
  2.2.3. Prove that $0_\F$, $1_\F$ and the additive inverse element $-a$ of any element $a \in F$ are unique.
}
\begin{solution}
  % ==== Place for solution here ====
\end{solution}


\problem{
  2.2.4. Prove $(-a)(-b) = ab$ for all $a, b \in \F$.
}
\begin{solution}
  % ==== Place for solution here ====
\end{solution}

\problem{
  2.2.5. Prove that the triple $(\Z_p, + \mod p, \cdot \mod p)$ is a field $\iff$ p is a prime. Use the fact that if two integers $m$, $n$ are relatively prime, then there
exist integers $r$ and $s$ such that $mr + ns = 1$.
}
\begin{solution}
  % ==== Place for solution here ====
\end{solution}


\problem{
  2.2.6. Prove that the 2-tuple $(S, \circ)$, where $S$ is the set of rotational symmetries of a square and $\circ$ is the composition operation, does indeed form a group. Moreover, find the symmetry
group of a rectangle.
}
\begin{solution}
  % ==== Place for solution here ====
\end{solution}


\problem{
  2.2.7. Characterise all idempotent elements of $(\cM_{2\times 2}(\R), +, \cdot)$.
}
\begin{solution}
  % ==== Place for solution here ====
\end{solution}


\problem{
  2.2.8. Prove that $\cC(\cR)$ is a subring of any ring $(\cR, +, \cdot)$.
}
\begin{solution}
  % ==== Place for solution here ====
\end{solution}


\problem{
  2.2.9. Let $(\cR, +, \cdot)$ be a ring and suppose that $\alpha \in \cR$ has a multiplicative inverse. Moreover,
suppose that $\alpha \in C(\cR)$. Is $\alpha^{-1} \in C(\cR)$?
}
\begin{solution}
  % ==== Place for solution here ====
\end{solution}


\problem{
  2.2.10. Does the set of functions with infinitely many roots form a subring of $\mathcal{F}(\R)$? Justify your answer.
}
\begin{solution}
  % ==== Place for solution here ====
\end{solution}


\problem{
  2.2.11. There is no function $g \in \text{supp}(\cF(\R))$ that acts as multiplicative identity for all elements in supp$(\cF(\R))$. Can you prove this?
}
\begin{solution}
  % ==== Place for solution here ====
\end{solution}


\problem{
  2.2.12. Determine the zero and the identity elements of the polynomial ring $(\cR[x], +, \cdot)$ where $\cR = \cC(\cM_{2\times 2}(\R))$.
}
\begin{solution}
  % ==== Place for solution here ====
\end{solution}

\problem{
  2.2.13. True/False. Justify:
  \begin{enumerate}
    \item Let $\F$ be any field. Suppose that $\alpha$ and $\beta$ are two elements of an extension field $\K$ of $\F$. Then $\F(\alpha)$ and
$\F(\beta)$ are always distinct extension fields of $\F$.
    \item Given any set $X$, there is no operation · that can be defined on $[X]^{\leq\left|X\right|}$ that makes it a group.
    \item Given a monoid $(\mathcal{S}, *)$ there exists a binary operation $\cdot$ on $\mathcal{S}$ such that $(\mathcal{S}, \cdot)$ no longer admits an identity element. If so, what assumption did you need to make?
    \item $(\R, \cdot)$ forms a group.
    \item $(\R, +)$ forms a group with the identity element being 1.
    \item Suppose $(\F, +, \cdot)$ is the field. Then $(\F, +)$ and $(\F, \cdot)$ both form a group.
  \end{enumerate}
}
\begin{solution}
  % ==== Place for solution here ====
\end{solution}

\problem{3.1.1. Had we defined vector addition by $\hat{A}\boxplus\hat{B}=[a_{ij}\cdot b_{ij}]$, i.e., pointwise multiplication of the entries in Example 3.1.3, would $(\mathcal{M}_{m\times n},\,\boxplus,\,\odot,\,(\mathbb{F},+,\cdot))$ form a vector space over $\mathbb{F}$?
}
\begin{solution}
  % ==== Place for solution here ====
\end{solution}


\problem{
3.1.2. Prove properties 2, 4 and 5 of Theorem 3.1.2, namely, for any vector space $V=(V,\boxplus,\odot,(\mathbb{F},+,\cdot))$ over field $\mathbb{F}$:
\begin{enumerate}
    \item[2.] The additive inverse of any $v\in V$ is unique.
    \item[4.] Any scalar of the zero vector yields the zero vector, i.e. $\forall\,\alpha\in\mathbb{F},\,a\odot\vec{0}=\vec{0}$.
    \item[5.] For all $v\in V$ and $c\in\mathbb{F}$, $(-c) \odot v = c \odot (-v) = -1 \odot (c \odot v)$.
\end{enumerate}
}

\begin{solution}
  % ==== Place for solution here ====
\end{solution}


\problem{
3.2.1. Prove Theorem 3.2.2, namely if $S\subseteq V$ then $\text{span}(S)\subseteq V$ is a subspace of $V$.
}
\begin{solution}
  % ==== Place for solution here ====
\end{solution}


\problem{
3.2.2. Prove Theorem 3.2.3, namely: Let $J$ be any index set (It is a set whose elements are used to enumerate the elements of another set.) and let $\{W_\alpha\mid\alpha\in J\}$ be a collection of subspaces of $V$. Then $\bigcap\limits_{\alpha\in J} W_{\alpha}$ is a subspace of $V$.
}
\begin{solution}
  % ==== Place for solution here ====
\end{solution}
\newpage

\problem{
3.2.3. Prove Theorem 3.2.4, namely: Let $W_1\subseteq W_2\subseteq W_3\subseteq\dots$ be an increasing chain of subspaces of $V$. Then $\bigcup\limits_{n\in\mathbb{N}}W_n$ is a subspace.
}
\begin{solution}
  % ==== Place for solution here ====
\end{solution}


\problem{3.3.1. Show that the set $\{e^x,\,e^{2x},\,e^{3x}\}$ is linearly independent in the vector space of $\mathcal{F}(\mathbb{R})$ over $\mathbb{R}$.}
\begin{solution}
  % ==== Place for solution here ====
\end{solution}


\problem{
3.4.1. Prove Proposition 3.4.1, namely, the following: \\
Let $T\colon V\to W$ be a linear transformation from $V$ to $W$. Then:
\begin{enumerate}
    \item For all $x,\,y\in V$, $T(x-y)=T(x)-T(y)$
    \item T maps the zero vector $\vec{0}_V$ of $V$ to the zero vector $\vec{0}_W$ of W.
    \item If $T'\colon W\to Z$ is also a linear transformation, their composition $T'\circ T$ is also a linear transformation from $V$ to $Z$ where $T'\circ T=T'(T(v))$. So for every vector $\vec{v} \in V$, we have $(T'\circ T)(\vec{v}):=T'(T(\vec{v}))$.
    \item If $V_1$ is a subspace of $V$ and $W_1$ is a subspace of $W$, then $T[V_1]$ is a subspace of $W$ and $T^{-1}[W_1]$ is a subspace of $V$.
\end{enumerate}
}
\begin{solution}
  % ==== Place for solution here ====
\end{solution}


\problem{
  3.4.2. Prove $$T: V \to W \text{ is injective } \iff \ker T = \{\vec{0}\}$$}
\begin{solution}
  % ==== Place for solution here ====
\end{solution}

\problem{
  3.4.3 $(\mathcal{L}(V,W), \boxplus, \odot, (\mathbb{F}, +, \cdot))$ forms a vector space over $\mathbb{F}$, where $\boxplus$ and $\odot$ are the pointwise addition of linear maps and scalar multiplication with the linear map in $W$, respectively. \\ \\
  That is given $T, S \in \mathcal{L}(V,W)$ and $\alpha \in \mathbb{F}$
    $$T \boxplus S := T \underbrace{+}_{W} S$$
    $$\alpha \odot T := \alpha \underbrace{\cdot}_{W} T$$
  So, given any $v \in V$ the image of $v$ under the linear maps $T \boxplus S$ and $\alpha \odot T$ would be $T(v) +_W S(v)$ and $\alpha \cdot_W T(v)$ respectively.\\ \\
  That is, $\cL(V,W)$ is closed under both $\boxplus$ and $\odot$. That is to show that $T\boxplus S$ and $\alpha \odot T$ are indeed linear transformations from V to W and all the axioms of a vector space are satisfied.
}
\begin{solution}
  % ==== Place for solution here ====
\end{solution}

\problem{
  3.4.4. Prove: \\ Let $V$ and $W$ be two vector spaces over $\F$ with $\text{dim}(V) = n$ and $\text{dim}(W ) = m$. Then, $\text{dim}(\cL(V, W )) = mn$\\\\
  Hint: Note that a linear tranformation $T$ is determined by its action on its basis elements of its domain. So, if $\mathcal{B}$ is a basis for $V$ over $\F$, and $T\in\cL(V, W)$, then as long as we know what $T(b)$ is for every $b \in \mathcal{B}$, (the basis could be finite or infinite) we could then determine the image of an arbitrary vector $v \in V$. Since $v = \sum\limits_{i=1}^{n} \alpha_i b_i$ for some $n \in w$, $b_i \in \mathcal{B}$, and $\alpha_i \in \mathbb{F}$, then by the properties of linear operators we have: $$T(v) = \sum_{i=1}^{n} \alpha_i T(b_i)$$
}
\begin{solution}
  % ==== Place for solution here ====
\end{solution}

\problem{3.4.5. Prove that a linear map $T\colon V\to W\text{ is invertible }\iff T\text{ is injective and surjective}$.}
\begin{solution}
  % ==== Place for solution here ====
\end{solution}

\problem{3.4.6. Let $V$ and $W$ be two vector spaces over a field $\F$ with dimensions $n$ and $m>3$ respectively. Construct three different linearly independent maps from $V$ to $W$.}
\begin{solution}
  % ==== Place for solution here ====
\end{solution}


\problem{3.4.7. Prove whether or not the triple $(\cL(\R,\,\R),\,+,\,\cdot)$ forms a subring of $\mathcal{F}(\R)$ under point-wise addition and multiplication.}
\begin{solution}
  % ==== Place for solution here ====
\end{solution}

\problem{
  3.6.1. Let $V$ and $W$ be two vector spaces over $\F$ with $\dim(V ) = n$ and $\dim(W ) = m$. Then
$\dim(\cL(V, W )) = mn$
}
\begin{solution}
  % ==== Place for solution here ====
\end{solution}

\problem{
  3.6.2. Suppose $A$ and $B$ are two similar matrices of size $n$. Prove whether or not they represent, or
can be viewed as matrix representations of the same linear transformation relative to suitable ordered bases.
}

\begin{solution}
  % ==== Place for solution here ====
\end{solution}

\problem{
  3.6.3. True/False. Justify
  \begin{enumerate}
    \item Let $\cB$ and $\cB'$ be two ordered bases for the vector space $\R^n$ over $\R$. Then $\det(\cM_{\cB}^{\cB'}) = 1$
    \item For all choices of ordered bases $\cB$ and $\cB'$ for $\R^n$, $\det(\cM_{\cB}^{\cB'})\neq 0$.
    \item Let $T \in \cL(V, W)$ be two vector spaces of dimension $n$ over a field $\F$. If $T$ is surjective, then it must be an
    isomorphism.
    \item $\mathbb{P}_{10}[x] \cong \R^{11}$.
    \item $T \in \cL(V, W)$ is an isomorphism $\iff \dim(V) = \dim(W)$.
    \item There are finite dimensional vector spaces $V$ and $W$ such that $\dim(\cL(V, W)) = 32$.
  \end{enumerate}
}
\begin{solution}
  % ==== Place for solution here ====
\end{solution}

\problem{
  4.2.1 For a basic vector space $V$ over $\C$, prove the following property of the inner product:
  $$\forall u, v, w \in V, \langle v + u, w \rangle = \langle v, w \rangle + \langle u, w \rangle$$
}
\begin{solution}
  % ==== Place for solution here ====
\end{solution}

\problem{
  4.2.2 Verify the following function is an inner product on $\cC_0([0,1])$
  $$\langle \cdot, \cdot \rangle: \cC_0([0,1]) \times \cC_0([0,1]) \to \R$$
  $$ \langle f(x), g(x) \rangle = \int_{0}^{1} f(x)g(x)dx $$
}
\begin{solution}
  % ==== Place for solution here ====
\end{solution}


\problem{
  4.2.3 Verify the following function is an inner product on $\C^n$
  $$ \vec{\alpha} = (\alpha_1, \alpha_2, \dots, \alpha_n), \quad \vec{\beta} = (\beta_1, \beta_2, \dots, \beta_n) \in \C^n$$
  $$\langle \vec{\alpha}, \vec{\beta} \rangle = \sum_{i=1}^{n} \overline{\alpha_i}{\beta_i}$$
}
\begin{solution}
  % ==== Place for solution here ====
\end{solution}

\problem{
  4.2.4 Could we have defined the inner product on $\C^n$ by $\langle \vec{\alpha}, \vec{\beta} \rangle = \sum_{i=1}^{n} {\alpha_i}\overline{\beta_i}$. Would it be an inner product over $\C^n$?
}
\begin{solution}
  % ==== Place for solution here ====
\end{solution}

\problem{
  4.2.5. Let $(V, \langle \cdot, \cdot, \rangle)$ be an inner product space. 
  Then V is a normed vector space with norm $\|v\| = \sqrt{\langle v, v \rangle}$. Prove the positive-definiteness of the norm.
}
\begin{solution}
  % ==== Place for solution here ====
\end{solution}

\problem{
  4.2.6. Could we ever obtain equality in the Cauchy-Schwarz inequality? i.e, $$|\langle u, v \rangle| \leq \sqrt{\langle u, u \rangle} \sqrt{\langle v, v \rangle} $$
}
\begin{solution}
  % ==== Place for solution here ====
\end{solution}

\problem{4.2.7. Prove the generalised version of the Pythagorean Theorem. i.e., Let $V$ be an inner product space. 
  Let $u, v$ be any two orthogonal vectors in $V$. Then:
  $$\|u_v\|^2 = \|u\|^2 + \|v\|^2$$ 
}
\begin{solution}
  % ==== Place for solution here ====
\end{solution}

\problem{
  4.2.8. Define the discrete metric $d$ on a vector space $V$ as follows:
  $$
      \forall u, v \in V, \quad d((u,v)) =
      \begin{cases}
      0 & \text{if } u = v, \\
      1 & \text{otherwise.}
      \end{cases}
  $$
  Verify that $d$ is indeed a metric on $V$.
}
\begin{solution}
  % ==== Place for solution here ====
\end{solution}

\problem{
  4.2.9. In Example 4.2.9, we define a metric on $\cC_0([0,1], \R)$ as follows:
  $$
  d(f(x),g(x)) \longmapsto \int_{0}^{1}|f(x)-g(x)| \psi(x)dx$$
  $$
  (\text{where } \psi: [0,1] \to (0, \infty) \text{ is any fixed continuous function and } \psi \neq 1 )
  $$
  Prove that function $d$ is indeed a metric on $\cC_0([0,1], \R)$. Make sure you understand it well. What happens if $\psi$ is not continuous?
}
\begin{solution}
  % ==== Place for solution here ====
\end{solution}

\problem{
  4.2.10. Let $(V, \|\cdot\|)$ be a normed vector space. Also, let $\|\cdot\|$ 
  follow the paralellogram law: 
  $$\| u + v \|^2 + \| u - v \|^2 = 2(\|u\|^2 + \|v\|^2)$$
  Define the inner product on $V$ as:
  $$(u,v) \longmapsto \langle u,v\rangle = \frac{1}{4}\left( \|u+v\|^2 - \|u-v\|^2 + i\|u-iv\|^2 - i\|u+iv\|^2 \right)$$
  Prove that $\langle \cdot, \cdot \rangle$ is indeed an inner product on $V$.
}

\end{document}